\subsection{Contextual straightening and every two-step path}

The selected column words in the source dictionary are
\[
\begin{array}{c|cc@{\qquad}c|cc}
\text{column}&V&W&\text{column}&V&W\\\hline
1&1&1&3&2&1\\4&2&2&12&11&21\\
13&21&11&23&21&12\\32&12&21&123&211&121\\
124&211&221&134&212&121\\234&212&221&&&
\end{array}
\]
These are diagram labels: they are not claims that the corresponding
nonstandard column equals just one canonical tensor word. For example
the column $(23)$ includes the source's additional $p^{-1}$ term.
The source's sign convention takes $(-p^{-1})^{\deg}$ times an
honest representative with no $3$--$2$ arc, or with its $3$--$2$
arc on the $W$ side, as positive. The displayed $T,T_1$ have degree
zero, while $T_2,Q$ have degree one. The diagram of $Q$ has its
$3$--$2$ arc on the $W$ side. Thus the displayed rescalings have
the source's positive orientation. The diagram classification places
$T,T_1,T_2,Q$ respectively in the pure $3$--$1$ $W$, $3$--$2$--$1$
$W$, no-external-arc, and pure $3$--$2$ $W$ cases after deleting
invariants; these are honest classes in the source basis.
Pair each $2$ with the next available $1$ to its right, recursively
matching $21$ pairs; equivalently scan left to right using a stack
of unmatched $2$ positions. The source's generator rule is
\[
 F_VD=\sum_{j=1}^{\varphi_V(D)}[j]\mathcal F_{(j)V}(D),
\]
where $\varphi_V$ is the number of unpaired $V$-ones. The $j$th term
changes the column containing that one by its column Kashiwara
operator, unless the changed word has an extra internal arc, in
which case the term is zero. The changes needed below are
$123\to134$, $124\to234$, $12\to32$, $1\to3$.
For an unpaired $(12)$ column, choosing the first of its two
$V$-ones creates an extra internal arc and gives zero; choosing
the second gives $(32)$.

For clarity we specify the contextual nonintegral rule used here.
For a consecutive pair of equal-height columns in one of the source
nonintegral strings, write the entire based product as $BCD$.
Let $I$ be the corresponding equal-height invariant, and let
$\omega$ be the unpaired word of $C$ on the indicated side. The
source's contextual projection identity, after taking the quotient,
is
\[
 BCD=-p^{-1}\begin{cases}B'ID&\omega=11,\\
 BID&\omega=12,\\BID'&\omega=22.\end{cases}
\]
In the $11$ instances below, $B'$ is obtained by changing the
unpaired $2$ in $B$ paired to the first local $1$ to a $1$; if
there is no such $2$, $B'=0$. The $12$ instance leaves both outside
words unchanged. These identities follow by projecting the source
formula $B\mathbin{\heartsuit_g}\widetilde C\mathbin{\heartsuit_g}D
=BCD+p^{-1}B'ID$ (and its $12$ version) to the quotient in which
the left-hand side is zero. This is the source corollary labelled
\texttt{c projecting to heartproj basis Kronecker}; the distinction
between the two based products in that formula is essential.

The pairs used here are $(134)(123)$, with $W$-word $11$ and
height-three invariant $(234)(123)$; $(134)(124)$ with $W$-word
$12$ and the same invariant; $(32)(12)$ with $V$-word $11$ and
height-two invariant $(34)(12)$; and $(3)(1)$ with $W$-word $11$
and height-one invariant $(4)(1)$. No isolated relation such as
$(3)(1)=0$ is substituted inside a longer based product.

The integral relations needed in the same quotient are
\begin{gather*}
 (234)(12)=-(124)(23),\qquad (134)(12)=-(124)(13),\\
 (32)(1)=(23)(1),\qquad (12)(23)=(23)(12).
\end{gather*}
They are the source's height-$(3,2)$, $(2,1)$ and $(2,2)$
degree-preserving relations. Their contextual products remain
relations: for integral relations the source's two based products
agree. Thus there is no suppressed nonintegral correction in these
four replacements.

In $T$, the prefix before its last two columns has exactly one
unpaired $W$-two, in column $(124)$. It is paired with the first
of the two suffix $W$-ones. The contextual rule therefore gives
\[
 (123)(123)(124)(12)(12)(3)(1)
 =-p^{-1}(123)(123)(123)(12)(12)(4)(1)=T_2.
\]
For $T_1$ the analogous unpaired $W$-two is in column $(23)$,
and changing it to a $1$ changes that column to $(13)$. Hence
\[
 (123)(123)(124)(23)(12)(3)(1)
 =-p^{-1}(123)(123)(124)(13)(12)(4)(1)=Q.
\]
This constructs the endpoint and its minus sign directly. The second
calculation is source case (F.1-1.5): its hypothesis holds because
$T_1$ has a $3$--$2$--$1$ $W$-arc and no pure $3$--$1$ $W$-arc.

We record how invariant counts exclude paths. An equal-height
invariant pair has all of its letters paired internally or within
the pair on both diagrams, with no exiting arc. An operator
$\mathcal F_{(j)V}$ changes an unpaired letter, so it cannot change
such a pair. Degree-preserving straightening moves preserve the
invariant record. In a nonintegral correction, the added pair is
an invariant and the contextual change is at an outside unpaired
letter; it therefore leaves pre-existing pairs intact. Integral
graded alternatives give zero. These facts prove, step by step,
that nonzero generator coefficients cannot decrease any of the
height-specific invariant counts. They also prove the direction
of containment used here independently of the reversed containment
sentence in source line 5007. In particular a vector with a
height-three or height-two invariant cannot reach $Q$, whose only
invariant is its final $(4)(1)$.

For $T$ the ordered unpaired $V$-one positions (column, word position)
are
\[
 (1,3),(2,3),(3,3),(4,1),(4,2),(5,1),(5,2),(6,1),(7,1).
\]
Put
\[
 H=-p^{-1}(123)(234)(123)(12)(12)(1)(1),\qquad
 R_0=(123)(123)(124)(12)(12)(1)(3).
\]
The following table lists all nine terms before multiplication by
their quantum integer $[j]$:
\[
\begin{array}{c|c|l}
j&\mathcal F_{(j)V}T&\text{verification}\\\hline
1&0&(134)(123),\ W:11,\ B'=0\\
2&H&(134)(124),\ W:12,\ \text{context unchanged}\\
3&-T_1&(234)(12)=-(124)(23)\\
4&0&\text{extra internal arc in column }4\\
5&0&(32)(12),\ V:11,\ B'=0\\
6&0&\text{extra internal arc in column }5\\
7&T_1&(32)(1)=(23)(1),\ \text{then commute }12,23\\
8&T_2&\text{the first contextual suffix calculation above}\\
9&R_0&\text{last unpaired one, the crystal term}
\end{array}
\]
All displayed zero-prefix assertions follow by the stack pairing:
for $j=1$ the prefix is empty; for $j=5$ the prefix has no unpaired
$V$-two. The vector $H$ has a height-three invariant, so none of its
nonzero successors can equal $Q$. This proves
$F_VT=([7]-[3])T_1+[8]T_2+[2]H+[9]R_0$, and only the last term
besides $T_1,T_2$ still requires a possible-return check.

There are eight unpaired $V$-ones in $R_0$, the first eight
positions in the preceding list; its last $3$ supplies an unpaired
$V$-two. Its $W$-diagram equals that of $T$. Put
\[
 H_R=-p^{-1}(123)(234)(123)(12)(12)(1)(3),\qquad
 R_1=(123)(123)(124)(23)(12)(1)(3),
\]
and $R_2=(123)(123)(124)(12)(12)(3)(3)$. Every term is
\[
\begin{array}{c|rrrrrrrr}
j&1&2&3&4&5&6&7&8\\\hline
\mathcal F_{(j)V}R_0&0&H_R&-R_1&0&0&0&R_1&R_2.
\end{array}
\]
The same local relations and zero-prefix calculations apply since
the last column lies outside them. The class $H_R$ has a height-three
invariant. The honest class $R_1$ has no invariant and one unpaired
$V$-two; $R_2$ has no invariant and two unpaired $V$-twos. The source
diagram classification identifies these honest classes, so neither
is $Q$, which has invariant record $(0,0,0,1)$ and no unpaired
$V$-two. This direct enumeration proves that the coefficient of
$Q$ in $F_VR_0$ is zero; no general monotonicity assertion about
unpaired twos is needed.

The ordered positions for $T_1$ are
$(1,3),(2,3),(3,3),(5,1),(5,2),(6,1),(7,1)$.
Its seven terms, with the source cases made explicit, are
\[
\begin{array}{c|l}
j&\text{outcome and reason it can or cannot equal }Q\\\hline
1&0:\ (134)(123),\ W:11,\ \text{empty prefix}\\
2&\text{a height-three invariant, case (F.3-3.1)}\\
3&\text{a height-two invariant, case (F.3-2)}\\
4&0:\ \text{extra internal arc in column }5\\
5&\text{a height-two invariant, case (F.2-1)}\\
6&Q:\ \text{the second contextual suffix calculation above}\\
7&R_1:\ \text{the crystal term, invariant count }i_1=0.
\end{array}
\]
The stated invariant cases can also be checked with the finite
relations: for $j=2$, apply the $W:12$ replacement
$(134)(124)\mapsto-p^{-1}(234)(123)$. For $j=3$, first use
$(234)(23)=-(234)(32)$ and then the contextual $V:11$ replacement
of $(32)(12)$, whose preceding unpaired $V$-two is in $(234)$;
it raises that column to $(124)$ and creates $(34)(12)$.
For $j=5$, use $(32)(1)=(23)(1)$ and then the $W:12$ replacement
$(23)(23)\mapsto-p^{-1}(24)(13)$, which creates a height-two
invariant. The extra integral relation here is the source's
$(3,2)$ bottom-weight relation, and the latter contextual relation
is its height-two nonintegral string. These computations account
for all indices and prove $a^{F_V}_{Q,T_1}=[6]$.

For $T_2$ the seven positions are
$(1,3),(2,3),(3,3),(4,1),(4,2),(5,1),(5,2)$; its existing
$(4)(1)$ invariant contributes none. Terms $j=1,2$ vanish because
their $(134)(123)$ $W:11$ replacements have no preceding unpaired
$W$-two. Terms $j=4,6$ vanish by extra internal arcs. Term $j=5$
vanishes by the $V:11$ $(32)(12)$ relation with no preceding
unpaired $V$-two. Term $j=7$ is the crystal vector
$-p^{-1}(123)(123)(123)(12)(32)(4)(1)$, with one unpaired
$V$-two; after removing its invariant its diagram has no external
arc, whereas $Q$ has a pure $3$--$2$ $W$-arc and no unpaired
$V$-two. Finally the exact $j=3$ calculation is
\[
 -p^{-1}(123)(123)(134)(12)(12)(4)(1)
 =p^{-1}(123)(123)(124)(13)(12)(4)(1)=-Q.
\]
Thus $a^{F_V}_{Q,T_2}=-[3]$. Every possible first step and every
remaining second step has now been accounted for. Therefore precisely
$T_1$ and $T_2$ contribute to the two-step path from $T$ to $Q$,
with four edge coefficients $[7]-[3]$, $[8]$, $[6]$, and $-[3]$.
